\documentclass[11pt,a4paper]{article}

% ============================================================
% PAPER 3 (draft/kerangka) — Bahasa Indonesia
% Self-Calibrating NIDS: Streaming Concept-Drift Detection +
%   Autonomous Few-Shot Adaptation untuk Trafik Cloud.
% Kelanjutan langsung Paper 1 (SFM + few-shot OFFLINE):
%   Paper 1 = adaptasi one-shot sebelum deploy.
%   Paper 3 = adaptasi ONLINE/continuous saat runtime (closed-loop).
% CATATAN: seluruh angka WAJIB dari eksperimen nyata; placeholder ditandai [hasil].
% ============================================================

\usepackage[utf8]{inputenc}
\usepackage[T1]{fontenc}
\usepackage[bahasa]{babel}
\usepackage{graphicx}
\usepackage{booktabs}
\usepackage{amsmath,amssymb}
\usepackage{siunitx}
\usepackage{hyperref}
\usepackage{geometry}
\usepackage{caption}
\usepackage{multirow}
\geometry{margin=2.5cm}
\hypersetup{colorlinks=true, linkcolor=blue, citecolor=blue, urlcolor=blue}

\title{\textbf{NIDS yang Mengalibrasi-Diri: Deteksi \emph{Concept Drift}
Streaming dan Adaptasi \emph{Few-Shot} Otonom untuk Trafik Cloud}}
% English working title: "Self-Calibrating NIDS: Streaming Concept-Drift
% Detection and Autonomous Few-Shot Adaptation for Cloud Traffic".

\author{Hero Yudo Martono, \quad Iwan Syarif, \quad Ferry Astika Saputra \\
\textit{Departemen Teknik Informatika, Politeknik Elektronika Negeri Surabaya}}
\date{Draft kerja --- \today}

\begin{document}
\maketitle

\begin{abstract}
\textbf{[KERANGKA — diisi setelah eksperimen.]} NIDS berbasis ML yang dikalibrasi
lintas-jaringan (mis. via \emph{Semantic Feature Mapping} dan \emph{few-shot},
\cite{paper1}) melakukan kalibrasi \emph{offline} sekali sebelum penyebaran. Namun
trafik cloud nyata mengalami \emph{concept drift} dinamis (pola siang/malam,
musiman), sehingga model statis lambat-laun menurun. Makalah ini mengusulkan
kerangka \emph{closed-loop}: detektor \emph{drift} \emph{streaming} berbasis
jendela-geser jarak Wasserstein ($W_1$) dan CUSUM pada aliran \emph{flow}
NFStream, yang memicu pembaruan \emph{micro-few-shot} otomatis dengan anggaran
label minimal (\emph{active learning}), disertai \emph{guardrail} anti-peracunan
untuk mencegah adaptasi keliru saat serangan berlangsung. [Ringkasan hasil
menyusul.]

\vspace{0.5em}
\noindent\textbf{Kata kunci:} NIDS, \emph{concept drift}, jarak Wasserstein,
\emph{online domain adaptation}, \emph{few-shot}, \emph{active learning}, cloud.
\end{abstract}

% ============================================================
\section{Pendahuluan}
\label{sec:intro}
\textbf{[Kerangka.]} Motivasi: kalibrasi \emph{offline} Paper~1~\cite{paper1}
tervalidasi memulihkan deteksi pada trafik AWS, tetapi \emph{one-shot}. Celah:
adaptasi berkelanjutan terhadap \emph{concept drift} runtime belum ditangani.

\paragraph{Gap yang ditutup.}
\begin{enumerate}
  \item \textbf{Drift tak-terdeteksi.} Model statis tidak tahu kapan
  distribusinya sudah bergeser jauh dari data latih.
  \item \textbf{Adaptasi manual/mahal.} Kalibrasi ulang biasanya butuh intervensi
  dan pelabelan besar.
  \item \textbf{Risiko adaptasi keliru.} Update saat serangan dapat
  ``mengajari'' model menerima serangan sebagai normal (\emph{poisoning}).
\end{enumerate}

\paragraph{Kontribusi (rencana).}
\begin{itemize}
  \item \textbf{Detektor \emph{drift} streaming} berbasis jendela-geser $W_1$ per
  fitur SFM $+$ CUSUM untuk drift lambat-kumulatif (\S\ref{sec:detector}).
  \item \textbf{Pemicu cerdas} yang membedakan \emph{benign drift} (mis.
  siang/malam) dari \emph{attack surge}, sehingga update hanya pada drift jinak
  (\S\ref{sec:trigger}).
  \item \textbf{Pembaruan \emph{micro-few-shot} otonom} dengan \emph{active
  learning} (label hanya untuk sampel paling informatif) $+$ \emph{guardrail}
  anti-\emph{poisoning} (\S\ref{sec:update}).
  \item \textbf{Evaluasi runtime multi-hari} di AWS: MCC terjaga seiring waktu
  vs baseline statis yang menurun (\S\ref{sec:eval}).
\end{itemize}

% ============================================================
\section{Karya Terkait}
\label{sec:related}
\textbf{[Kerangka.]} (i) \emph{Concept drift} pada \emph{data stream} (ADWIN,
DDM, Page-Hinkley); (ii) deteksi \emph{drift} berbasis jarak distribusi
($W_1$/MMD); (iii) \emph{online}/\emph{continual} learning pada NIDS;
(iv) \emph{active learning} untuk NIDS berlabel-terbatas; (v) posisi kami:
menutup \emph{loop} deteksi$\rightarrow$adaptasi dengan biaya label minimal dan
guardrail keamanan, di atas fondasi SFM$+$few-shot Paper~1.

% ============================================================
\section{Metode}
\label{sec:method}

\subsection{Detektor Drift Streaming ($W_1$ + CUSUM)}
\label{sec:detector}
\textbf{[Kerangka + persamaan.]} Untuk tiap fitur SFM $j$, hitung $W_1$ antara
distribusi jendela-geser terkini $\mathcal{W}_t$ (ukuran $w$) dan distribusi
acuan latih $\mathcal{D}_{\text{ref}}$:
\begin{equation}
\label{eq:driftscore}
S_t \;=\; \frac{1}{d}\sum_{j=1}^{d} W_1\!\big(\mathcal{W}_t^{(j)}, \mathcal{D}_{\text{ref}}^{(j)}\big).
\end{equation}
CUSUM mengakumulasi deviasi $S_t$ terhadap garis dasar untuk menangkap drift
lambat; alarm dipicu bila statistik kumulatif melewati ambang $h$. (Detail
kalibrasi $w$, $h$, dan garis dasar menyusul dari eksperimen notebook \texttt{30\_}.)

\subsection{Pemicu Cerdas: Benign Drift vs Attack Surge}
\label{sec:trigger}
\textbf{[Kerangka.]} Aturan pembeda (hipotesis, akan divalidasi): \emph{benign
drift} = pergeseran perlahan pada fitur temporal/laju \emph{tanpa} lonjakan laju
alarm; \emph{attack surge} = pergeseran mendadak berbarengan lonjakan alarm.
Hanya \emph{benign drift} yang memicu adaptasi; \emph{attack surge} memicu
mitigasi, bukan pembaruan model.

\subsection{Pembaruan Micro-Few-Shot Otonom + Guardrail}
\label{sec:update}
\textbf{[Kerangka.]} Saat pemicu jinak aktif: (i) \emph{active learning} memilih
$k$ sampel paling tak-pasti (dekat batas keputusan) untuk dilabeli (anggaran
kecil); (ii) model diperbarui \emph{inkremental} (few-shot) memakai resep Paper~1;
(iii) \emph{guardrail}: model kandidat divalidasi pada \emph{holdout} tepercaya
sebelum dipromosikan, mencegah \emph{poisoning}.

% ============================================================
\section{Eksperimen dan Hasil}
\label{sec:eval}
\textbf{[Kerangka — angka menyusul dari notebook \texttt{30\_}++ dan AWS multi-hari.]}
\begin{itemize}
  \item \textbf{PoC deteksi drift (offline, murah):} simulasikan \emph{stream}
  berurutan CIC$\rightarrow$UNS$\rightarrow$AWS; tunjukkan $S_t$ (Persamaan~\ref{eq:driftscore})
  melonjak di titik transisi domain. (Notebook \texttt{30\_drift\_detector\_poc}.)
  \item \textbf{Simulasi closed-loop:} bandingkan MCC model statis vs
  auto-adaptif sepanjang \emph{stream}.
  \item \textbf{Validasi AWS multi-hari:} capture berkelanjutan (mis. 3--7 hari),
  ukur MCC/FAR seiring waktu; buktikan auto-adapt menahan penurunan.
\end{itemize}

% ============================================================
\section{Diskusi \& Keterbatasan}
\label{sec:discussion}
\textbf{[Kerangka.]} Bahas: biaya label \emph{active learning}, risiko
\emph{poisoning} dan efektivitas guardrail, stabilitas \emph{loop} (hindari
osilasi update), serta batas generalisasi ke topologi lain.

% ============================================================
\section{Kesimpulan}
\label{sec:conclusion}
\textbf{[Kerangka.]} Ringkas kontribusi closed-loop dan hasil utama setelah
eksperimen selesai.

% ============================================================
\begin{thebibliography}{9}
\bibitem{paper1} H.~Y.~Martono, I.~Syarif, and F.~A.~Saputra, ``Generalisasi
Cross-Dataset pada NIDS Menggunakan Semantic Feature Mapping dan Few-Shot
Adaptation dengan Validasi Trafik Cloud Nyata,'' \emph{naskah kerja} (Paper~1),
2026.
\bibitem{layeghy} S.~Layeghy, M.~Gallagher, and M.~Portmann, ``Benchmarking the
benchmark --- Comparing synthetic and real-world network IDS datasets,''
\emph{J. Inf. Secur. Appl.}, vol.~80, art.~103689, 2024.
\bibitem{bendavid} S.~Ben-David \emph{et al.}, ``A theory of learning from
different domains,'' \emph{Mach. Learn.}, vol.~79, pp.~151--175, 2010.
% TODO: tambah rujukan concept drift (ADWIN, DDM, Page-Hinkley), MMD, active learning.
\end{thebibliography}

\end{document}
